
\documentclass[11pt,a4paper,onecolumn]{article}

% ============================================================
% Page Layout and Typography
% ============================================================
\usepackage[margin=1in]{geometry}
\usepackage[T1]{fontenc}
\usepackage[utf8]{inputenc}
\usepackage{lmodern}
\usepackage{microtype}

% ============================================================
% Mathematics
% ============================================================
\usepackage{amsmath}
\usepackage{amssymb}
\usepackage{amsthm}
\usepackage{mathtools}
\usepackage{bm}

% ============================================================
% Algorithms and Pseudocode
% ============================================================
\usepackage{algorithm}
\usepackage{algpseudocode}

% ============================================================
% Figures and Tables
% ============================================================
\usepackage{graphicx}
\usepackage{float}
\usepackage{booktabs}
\usepackage{caption}
\usepackage{subcaption}
\usepackage{tikz}

% ============================================================
% Lists and Formatting
% ============================================================
\usepackage{enumitem}
\usepackage{xcolor}

% ============================================================
% References and Hyperlinks
% ============================================================
\usepackage[numbers,sort&compress]{natbib}
\usepackage[
    colorlinks=true,
    linkcolor=blue,
    citecolor=blue,
    urlcolor=blue
]{hyperref}
\usepackage[capitalise,noabbrev]{cleveref}

% ============================================================
% Theorem Environments
% ============================================================
\theoremstyle{plain}
\newtheorem{theorem}{Theorem}[section]
\newtheorem{lemma}[theorem]{Lemma}
\newtheorem{proposition}[theorem]{Proposition}
\newtheorem{corollary}[theorem]{Corollary}

\theoremstyle{definition}
\newtheorem{definition}[theorem]{Definition}
\newtheorem{example}[theorem]{Example}

\theoremstyle{remark}
\newtheorem{remark}[theorem]{Remark}
\newtheorem{observation}[theorem]{Observation}

% ============================================================
% Bin Packing Notation
% ============================================================
\newcommand{\eps}{\varepsilon}
\newcommand{\I}{\mathcal{I}}

\DeclareMathOperator{\OPT}{OPT}
\DeclareMathOperator{\ALG}{ALG}
\DeclareMathOperator{\SIZE}{SIZE}

% Common number sets
\newcommand{\R}{\mathbb{R}}
\newcommand{\N}{\mathbb{N}}
\newcommand{\Z}{\mathbb{Z}}

% ============================================================
% Title Information
% ============================================================
\title{
    \textbf{CS4234 Mini-Project: Approximation Algorithms for Bin Packing}
}

\author{
    Author 1 \and
    Author 2 \and
    Author 3 \and
    Author 4
}

\date{November 2026}

\begin{document}

\maketitle

\section{Introduction}

The \textbf{Bin Packing Problem} is a classical combinatorial optimization problem with applications in resource allocation, scheduling, and storage management. Given a collection of items with sizes \(s_1,s_2,\ldots,s_n\in(0,1]\), the objective is to partition them into the minimum number of bins, each having unit capacity. Although the problem has a simple formulation, finding an optimal packing is NP-hard. This motivates the study of efficient approximation algorithms that produce solutions close to the optimum.

A central challenge in designing approximation algorithms for bin packing is balancing computational efficiency with solution quality. While simple greedy algorithms can construct feasible packings efficiently, they do not necessarily achieve solutions arbitrarily close to the optimum. An important line of research therefore studies *asymptotic approximation guarantees*, which permit an additive error while seeking a multiplicative approximation ratio approaching one.

A major breakthrough was achieved by Fernández de la Vega and Lueker (1981) \cite{fernandezdelavegaBinPackingCan1981}, who developed an asymptotic polynomial-time approximation scheme for bin packing. Their approach separates large and small items, applies *linear grouping* to reduce the number of distinct item sizes, and rounds item sizes to obtain a simplified instance. This reduction makes it possible to enumerate feasible bin configurations and solve the resulting packing problem efficiently. For any fixed \(\varepsilon>0\), their algorithm achieves a packing using at most
$ (1+\varepsilon)\operatorname{OPT}(I)+O(\varepsilon^{-2}) $ bins, with running time linear in the number of items when \(\varepsilon\) is fixed.

Subsequently, Karmarkar and Karp (1982)\cite{karmarkarEfficientApproximationScheme1982} developed a stronger approximation algorithm based on linear programming relaxation and iterative rounding techniques. By exploiting the structure of fractional packing solutions and carefully handling residual instances, their algorithm achieves the improved additive guarantee $\operatorname{OPT}(I)+O(\log^2\operatorname{OPT}(I))$. This result demonstrates how linear programming can be used to obtain substantially stronger approximation guarantees.

In this report, we study the fundamental techniques underlying these two algorithms. We first examine the de la Vega–Lueker algorithm, beginning with the restricted bin packing problem and proceeding to item grouping, rounding, and the treatment of small items. We then investigate the main ideas behind the Karmarkar–Karp algorithm, particularly its configuration-based linear programming formulation and the handling of residual packings. Throughout the report, we emphasize intuitive explanations, illustrative examples, mathematical justification of approximation guarantees, and analysis of running-time complexity. Our goal is to clarify how the two algorithms progressively improve the approximation of the bin packing problem.


% \begin{definition}[Bin Packing]
% Given an instance $\I = \{a_1,\ldots,a_n\}$,
% where $a_i \in (0,1]$, the bin packing problem
% asks for a partition of the items into the minimum
% number of bins such that the total size of items
% in each bin is at most $1$.
% \end{definition}

% \begin{theorem}
% For every $\eps > 0$, there exists an algorithm
% with an asymptotic approximation guarantee
% of $1+\eps$ for the bin packing problem.
% \end{theorem}

% \begin{proof}
% The proof proceeds by ...
% \end{proof}


\section{Preliminaries}

\subsection{Notations}

\subsection{Greedy Proof}

\subsection{Relaxed Bin Packing}

\section{De la Vega--Lueker Algorithm}
\cite{fernandezdelavegaBinPackingCan1981}

\section{Karmarkar--Karp Algorithm}
\cite{karmarkarEfficientApproximationScheme1982}

\section{Conclusion}

% ============================================================
% References
% ============================================================
\bibliographystyle{plainnat}
\bibliography{reference}

% ============================================================
% Appendix (outside the main-body page limit)
% ============================================================
\appendix

\section{Supplementary Proofs}

% ============================================================
% AI Disclosure (outside the 10-page limit)
% ============================================================
\section*{AI Disclosure}

% Describe how AI tools were used in this project.

\end{document}

